% Part of main.tex; compile main.tex.
% ---------------------------------------------------------------------------
\section{Evaluation}
\label{sec:eval}

\notebox{Scope}{All results in Section~\ref{sec:results} are measured on a
single node, with one run per configuration. PDC is measured with one
server in the main, ratio, size, and GPU experiments, but not at 48\,GiB per
step, not at compute-to-I/O ratios above 1.25, and not in the core-allocation
experiment (beyond the SMT-sibling placement of its clients). Its asynchronous
run with the halo exchange did not complete (Section~\ref{sec:impl}).}

In this section, we describe our research questions, the machine and
software used in our experiments, and the asynchronous I/O implementations
we evaluate. Sections~\ref{sec:surrogates} to~\ref{sec:metrics} then
describe the compute surrogates, the workload, and the metrics, and
Section~\ref{sec:results} presents the results.

\subsection{Research questions}
Our evaluation is guided by four research questions (RQs):
\begin{itemize}
  \item \textbf{RQ1:} How much longer does an application run than a
        sleep-based evaluation predicts, when its computation consumes CPU
        time, memory bandwidth, or network bandwidth?
  \item \textbf{RQ2:} Where does the cost of interference appear, as a
        slower computation or as exposed I/O, and which resource is
        responsible for each asynchronous I/O implementation?
  \item \textbf{RQ3:} How do conclusions drawn from sleep-based kernels
        compare with those from a mini-application that combines
        computation, memory traffic, and communication?
  \item \textbf{RQ4:} How sensitive are the results to the core allocation
        and to whether the computation runs on the CPU or the GPU?
\end{itemize}

\subsection{Machine setup}
\label{sec:machine}
All experiments run on a single dual-socket node. The hardware is described
in Table~\ref{tab:hw}, and the software versions and configuration are
described in Table~\ref{tab:sw}.

\begin{table}[t]
\centering
\small
\setlength{\tabcolsep}{3pt}
\begin{tabular}{@{}p{0.27\columnwidth}p{0.68\columnwidth}@{}}
\toprule
Component & Configuration \\
\midrule
CPU & 2$\times$ Intel Xeon Gold 5418Y (Sapphire Rapids), 24 cores each,
      2.0\,GHz base; 48 cores, 96 hardware threads (SMT on) \\
Caches & 48\,KiB L1d and 2\,MiB L2 per core; 45\,MiB L3 per socket \\
Memory & 256\,GB, two NUMA domains (128\,GB each); inter-node distance 21 \\
GPU & NVIDIA RTX 5000 Ada Generation, 32\,GB, PCIe Gen4 $\times$16 \\
Storage & 4$\times$ SanDisk PC SN8000S 4\,TB NVMe SSDs (PCIe Gen4 $\times$4 each)
          in Linux MD RAID-0 (512\,KiB chunk), 16\,TB \\
File system & ext4 (\texttt{stripe=512}, \texttt{relatime}) \\
Page cache & bounded per run by a cgroup~v2 memory limit (footprint + 24\,GiB);
             \texttt{vm.dirty\_ratio}=20, \texttt{dirty\_background\_ratio}=10 \\
Other & transparent huge pages \texttt{madvise}; no other user jobs during runs \\
\bottomrule
\end{tabular}
\caption{Hardware configuration of the test node.}
\label{tab:hw}
\end{table}

\begin{table*}[t]
\centering
\small
\setlength{\tabcolsep}{4pt}
\begin{tabular}{@{}l l p{0.33\textwidth} p{0.33\textwidth}@{}}
\toprule
Software & Version & Build configuration & Runtime configuration \\
\midrule
Ubuntu / Linux & 26.04 LTS / 7.0.0-29 & distribution kernel & THP \texttt{madvise}; cgroup~v2 memory limit per run \\
glibc & 2.43 & distribution & POSIX AIO: 1 helper thread (\texttt{aio\_init}) \\
GCC & 15.2.0 & distribution & \texttt{-O2 -march=native} for the benchmark \\
MPICH & 4.3.2 & distribution, \texttt{ch4:ucx}, Hydra & \texttt{MPI\_THREAD\_MULTIPLE}; explicit per-rank binding; ROMIO \texttt{cb\_config\_list=*:8}, \texttt{cb\_nodes=8}; async progress thread for non-blocking MPI-IO \\
UCX / hwloc & 1.20.0 / 2.13.0 & distribution & defaults \\
HDF5 & 1.14.6 & parallel, thread-safe, no compression, Release & MPI-IO VFD; collective data and metadata writes \\
HDF5 async VOL & git \texttt{90fbba9} (2026-02-10) & \texttt{ENABLE\_WRITE\_MEMCPY=ON}; wait back-off patch (Sec.~\ref{sec:impl}) & \texttt{HDF5\_ASYNC\_EXE\_DCLOSE=1}; event sets, one step in flight \\
Argobots & 1.2 & default & one execution stream per process (async VOL) \\
ADIOS2 & 2.12.1 & MPI on; SST, Python, Fortran, compression off; opt-in fsync patch (inactive) & BP5; \texttt{AggregationType=TwoLevelShm}, \texttt{NumAggregators=8}; \texttt{AsyncWrite=Guided} (async) \\
PDC & 0.7-398 (git \texttt{1ad9b312}) & MPI on; Mercury 2.2.0, libfabric 1.18.0 & \texttt{ofi+tcp}; 1 server on a reserved core; cache 24\,GiB \\
OpenBLAS & 0.3.34 & \texttt{TARGET=SAPPHIRERAPIDS}, \texttt{USE\_THREAD=0} & single-threaded DGEMM \\
CUDA & 13.3 (nvcc 13.3.73) & \texttt{-arch=native} (sm\_89) & driver 595.91.07; spin or blocking-sync scheduling \\
Python & 3.14.4 & NumPy 2.5.2, Matplotlib 3.11.1 & analysis and plotting \\
\bottomrule
\end{tabular}
\caption{Software configuration: versions, build options, and runtime settings
used in all experiments. Build options are set in \srcBuildHdf (HDF5),
\srcBuildVol (async VOL), and \srcBuildAdios (ADIOS2), and runtime settings in
\srcMatrixEnv and \srcAsyncEnv.}
\label{tab:sw}
\end{table*}

\paragraph{Process placement.} Unless stated otherwise, we run 48 MPI ranks.
Each rank is bound to one hardware thread on a distinct physical core, and
the ranks alternate between the two sockets. Background I/O threads inherit
the binding of their rank, and therefore share its hardware thread. We call
this configuration \emph{fully subscribed}. For RQ4, we also evaluate two
other placements. In the \emph{SMT sibling} placement, each rank is bound to
both hardware threads of its core. In the \emph{core reserved} placement,
we run 24 ranks, each bound to two physical cores on the same socket. We
also ran an unpinned placement, in which the operating system decides where
each thread runs, but we exclude it. Without pinning, threads were placed
differently during calibration and during the measurement, so the compute
phases reached only 30\% to 75\% of their target length. The GPU experiments use 4 ranks that share the single GPU. The binding
lists are constructed in \srcBinding.

\paragraph{Storage behavior.} Due to the 256\,GB of memory on our node, the
page cache could absorb most writes and turn the I/O into a memory copy. In
contrast, the client of a parallel file system can cache only a small part
of a checkpoint. For this reason, we run each experiment in a cgroup~v2
scope whose memory limit is the memory footprint of the benchmark plus
24\,GiB. Because page-cache pages are charged to the cgroup, this limit
bounds the cache available to a run. Once the cache is full, writers stall
in memory reclaim until dirty pages have been written back to the RAID
array. It is important to note that dirty-page throttling itself remains
global, at 20\% and 10\% of the memory, because the \texttt{io} controller is
not delegated to user cgroups on this system. In addition, no backend calls
\texttt{fsync} after each step. The HDF5 asynchronous VOL executes file
flushes synchronously by design, so per-step durability would make it
blocking. Every run starts with an empty output directory after a
\texttt{sync}, as shown in \srcCgroup, and the memory limit is computed in
\srcMemMax. As a result, small outputs of up to tens of gigabytes per run
complete largely at memory speed, as they would with a burst buffer, while
larger outputs are limited by write-back. The size sweep in
Section~\ref{sec:size} covers both regimes.

\subsection{Asynchronous I/O implementations}
\label{sec:impl}
Each backend has a blocking and an asynchronous mode behind a common
interface, \texttt{io\_backend}, shown in \srcHeaderBackend. All runs
initialize MPI with \texttt{MPI\_THREAD\_MULTIPLE} (\srcThreadInit). We
describe each backend below.
\begin{itemize}
  \item \textbf{POSIX.} The blocking mode writes with \texttt{pwrite} to a
        shared file per step, at the offset of the rank. The asynchronous
        mode copies the buffer to a staging buffer and writes it with POSIX
        AIO (\texttt{aio\_write}), which is executed by a glibc helper
        thread. The write path is shown in \srcPosixIssue, and the
        completion in \srcPosixWait.
  \item \textbf{MPI-IO.} The blocking mode writes one shared file per step
        with the collective \texttt{MPI\_File\_write\_at\_all}. The
        asynchronous mode writes a staged copy of the buffer with the
        non-blocking \texttt{MPI\_File\_iwrite\_at\_all} (\srcMpiioIssue).
        ROMIO uses collective buffering with 8 aggregators on the node
        (\texttt{cb\_config\_list=*:8}), which matches the ADIOS2 setting.
        Because a non-blocking collective write only advances inside MPI
        calls, the asynchronous mode enables the progress thread of MPICH
        (\texttt{MPIR\_CVAR\_ASYNC\_PROGRESS=1}). Similar to the other
        background engines, this thread inherits the binding of its rank.
  \item \textbf{HDF5.} Both modes write one shared file with a new dataset
        per step, using collective \texttt{H5Dwrite} and collective
        metadata. The asynchronous mode uses \texttt{H5Dwrite\_async} with
        event sets, which are executed by one Argobots background thread per
        process with write-buffer copying enabled~\cite{tang2022async}.
        Following the guidance of the connector for periodic
        output~\cite{volasync_docs}, we set
        \texttt{HDF5\_ASYNC\_EXE\_DCLOSE=1}, so that background execution
        starts when the dataset is closed instead of after an idle-detection
        period. The write path is shown in \srcHdfIssue, the wait in
        \srcHdfWait, and the connector selection in \srcAsyncEnv.
  \item \textbf{ADIOS2.} The blocking mode uses the BP5 engine with
        \texttt{AsyncWrite} disabled, and the asynchronous mode uses
        \texttt{AsyncWrite=Guided}, in which the compute phase is bracketed
        with computation-block hints~\cite{godoy2020adios2}. Both modes use
        two-level shared-memory aggregation with 8 aggregators. With the
        default of one aggregator per node, a single process writes all of
        the data of the node, which limits BP5 to about a third of the
        bandwidth of POSIX on our single node. The engine setup is shown in
        \srcBpInit, the per-step path in \srcBpIssue, the hints in
        \srcBpHints, and the aggregation settings in \srcMatrixEnv.
  \item \textbf{PDC.} The PDC backend uses synchronous or asynchronous
        region transfers, with one PDC server per node on a reserved
        core~\cite{tang2018pdc} (\srcPdcIssue). Following the PDC
        deployment model, we run 1 server on a reserved core and 47 client
        ranks on the remaining cores. The server cache is capped at
        24\,GiB, the same budget as the page cache of the other backends,
        and the cache is flushed before the timed run ends. Each client
        rank runs a Mercury progress thread during asynchronous transfers,
        which polls the network with a 1\,ms pause between polls. The
        asynchronous run with the halo exchange hung in two attempts and
        was stopped after 30 minutes each time (\srcRunPdc), so PDC has no
        halo result.
\end{itemize}

\paragraph{A wait-path defect in the HDF5 asynchronous VOL.} In early
runs, asynchronous HDF5 was slower than synchronous HDF5 even with no
computation, and the waits were quantized in multiples of one second. By
tracing system calls, we found that \texttt{H5VL\_async\_request\_wait},
which is invoked by \texttt{H5ESwait}, sleeps for a full second
(\texttt{usleep(1000000)}) whenever it fails to reacquire the global lock of
HDF5 from the background thread. In addition, it polls the task status every
100\,ms. This penalty applies only when the application waits while
background I/O is still running. Because this happens more often under
contention than with \texttt{sleep}, the unpatched connector would
exaggerate the effect we study. We therefore reduce both intervals to 1\,ms.
With this change, asynchronous HDF5 without computation performs on par
with synchronous HDF5. The patch is shown in \srcVolPatch, and it is applied
by the build script (\srcBuildVol).
